\documentclass[15pt]{scrartcl}
\usepackage[sexy]{evan}
\newcommand{\R}{\mathbb{R}}
\newcommand{\N}{\mathbb{N}}
\newcommand{\Q}{\mathbb{Q}}
\usepackage{relsize}
\begin{document}
\title{Hilbert's Theorem 90}
\date{\today}
\author{Jiawen Huang}
\maketitle

\section{Why do we care about Hilbert's Theorem 90?}
\begin{enumerate}
    \item What is the norm of an element in $\Q(\sqrt{2})$? Easy.
    \item But what is the norm of an element in $\Q(\sqrt[3]{2})$? How do we define it in general? Well, $N(a+b\sqrt[3]{2}+c\sqrt[3]{4})=a^3+2b^3+2c^3-6abc$. Not so easy. Why is that? 
    \item Find all the rational solutions to the equation $x^2+dy^2=1$ for \textbf{ANY} integer $d$.
    \item Classify all the elements of norm $1$ over Cyclic Galois extension.
\end{enumerate}
\section{Prerequisites}
\begin{enumerate}
    \item Galois Theory
    \item Separable Extensions 
\end{enumerate}
\section{What is Norm?}
\begin{enumerate}
    \item What is the norm of an element in $\Q(\sqrt{2})$?  
    \item But what is the norm of an element in $\Q(\sqrt[3]{2})$? How do we define it in general?
    \item So our definition should meet the basic property: multiplicative: $N(\alpha\beta) = N(\alpha)N(\beta)$
    \item In fact, for any finite extension $k\subseteq F$, the norm of $\alpha \in F$ over $k$ is defined as follows:
    \\
    We treat $\alpha$ as a $k$-linear transformation from $F$ to $F$ by multiplication, i.e. $f_\alpha: F \to F$ defined by $f_\alpha(x) = \alpha x$. Since $F$ is a vector space over $k$ of finite dimension, we can consider the determinant of this linear transformation. This determinant is defined to be the norm of $\alpha$ over $k$, denoted as $N_{F/k}(\alpha)$.
    \item So how do we find the determinant of this linear transformation? One solution is to look at the characteristic polynomial of this linear transformation. 
    \begin{example}
    Find the norm of $a+b\sqrt{2}$ in $\Q(\sqrt{2})$ over $\Q$.
    \\

    The norm of $a+b\sqrt{2}$ in $\Q(\sqrt{2})$ over $\Q$ is $a^2-2b^2$. This is because the characteristic polynomial of $a+b\sqrt{2}$ is the same as the minimal polynomial of $a+b\sqrt{2}$ over $\Q$ (degree considerations), which is $x^2-2ax+(a^2-2b^2)$. 
    \end{example} 
    \begin{example}
        Find the norm of $t=a+b\sqrt[3]{2}+c\sqrt[3]{4}$ in $\Q(\sqrt[3]{2})$ over $\Q$. Let's compute the characteristic polynomial. We have
\[\begin{pmatrix}
a & 2c & 2b \\
b & a & 2c \\
c & b & a
\end{pmatrix}\]
So the characteristic polynomial is $x^3-3ax^2+3(a^2-bc)x-(a^3+2b^3+2c^3-6abc)$. So the norm is $a^3+2b^3+2c^3-6abc$.
    \end{example}
    \begin{example}
        Find the norm of $a+b i$ in $\Q(i)$ over $\Q$. The characteristic polynomial is $x^2-2ax+(a^2+b^2)$. So the norm is $a^2+b^2$.
    \end{example}
    \begin{example}
        Find the norm of an element in $\Q(\sqrt{2},i)$ over $\Q$. Note that $\Q(\sqrt{2},i)$ is a degree 4 extension over $\Q$. A basis is given by $\{1,\sqrt{2},i,i\sqrt{2}\}$. So the characteristic polynomial of $a+b\sqrt{2}+ci+di\sqrt{2}$ is given by the determinant of the following matrix:
        \[\begin{pmatrix}
a & 2b & -c & -2d \\
b & a & -d & -c \\
c & 2d & a & 2b \\
d & c & b & a
\end{pmatrix}\]
So the norm is $a^4+4b^4+c^4+4d^4-8a^2b^2+2a^2c^2+8b^2c^2-8a^2d^2+8b^2d^2+2c^2d^2$. 
\\

But what if I am only interested in the norm of $a+ci$? When $b=d=0$, the matrix simplifies to
\[\begin{pmatrix}
a & 0 & -c & 0 \\
0 & a & 0 & -c \\
c & 0 & a & 0 \\
0 & c & 0 & a
\end{pmatrix}\]
So the norm is $(a^2+c^2)^2$.
\\
Wait a minute. This is very similar to the norm of $a+ci$ in $\Q(i)$ over $\Q$. Is that a coincidence? Is there a faster way to compute it?
\\

Well the answer is YES. 
    \end{example}
\begin{lemma}
    If $k\subseteq k(\alpha) \subseteq F$ are finite extensions, and $[F:k(\alpha)]=r$, then for any $\beta \in k(\alpha)$, we have
    \[
    N_{F/k}(\beta) = N_{k(\alpha)/k}(\beta)^r.
    \]
\end{lemma}
The key of the proof is to write down a basis of $F$ over $k$ using a basis of $k(\alpha)$ over $k$ and a basis of $F$ over $k(\alpha)$.
    \item With the Lemma, we don't have to compute the matrix and can use the norm in smaller extensions to find the norm in larger extensions. To be specific, we can compute the minimal polynomial of $\alpha$ over $k$. The constant term of this minimal polynomial (up to a sign) gives the norm of $\alpha$ over $k\subseteq k(\alpha)$.
    \\

    Then, we use the lemma, $N_{F/k}(\alpha) = N_{k(\alpha)/k}(\alpha)^{[F:k(\alpha)]}$ to get the norm of $\alpha$ over $k\subseteq F$.
    \begin{example}
        Find the norm of $a+ci$ in $\Q(\sqrt{2},i)$ over $\Q$.
        \\

        The minimal polynomial of $a+ci$ over $\Q$ is $x^2-2ax+(a^2+c^2)$. So the norm of $a+ci$ over $\Q(i)$ is $a^2+c^2$. Since $[\Q(\sqrt{2},i):\Q(i)]=2$, by the lemma, the norm of $a+ci$ over $\Q$ is $(a^2+c^2)^2$.
    \end{example}
\end{enumerate}
\section{When Separability comes into play}

\end{document}